Fix \(\chi\). Along \(0\le\epsilon\le10^{-3}\), the source probability is \(1/2\), the observational treatment kernel is bounded below by \(3/8\), and every proxy kernel appearing in either design is bounded away from zero by a positive constant depending only on \(\chi\). The two designs have finitely many cells. Hence there is \(c_\chi>0\) such that every observational cell of \((A,X,Z_{M_d})\), for both \(d\), has probability at least \(c_\chi\), uniformly along the triangular sequence.

The training and validation observational counts are both proportional to \(m\) in probability. A binomial concentration bound and a finite union bound therefore show that every training cell count is proportional to \(m\) with probability tending to one. Since \(|Y|=1\), each saturated training-cell mean satisfies, uniformly over the finitely many cells and designs,
\[
 \widehat f_{d,m}^A(v)-E(Y\mid (A,X,Z_{M_d})=v,G=O)=O_P(m^{-1/2}). \tag{1}
\]
Condition on the training half. The held-out squared losses are independent, bounded by four, and their number is proportional to \(m\). Their empirical mean differs from its conditional expectation by \(O_P(m^{-1/2})\), uniformly over the two designs. If \(f_d^A=E(Y\mid A,X,Z_{M_d},G=O)\), the conditional-expectation projection identity gives
\[
 E\{(Y-\widehat f_{d,m}^A)^2\mid G=O,\mathcal I_{m,O}^{\mathrm{tr}}\}
 -R_d^A(P_{\chi,\epsilon_m})
 =E\{(\widehat f_{d,m}^A-f_d^A)^2\mid G=O,\mathcal I_{m,O}^{\mathrm{tr}}\}
 =O_P(m^{-1}). \tag{2}
\]
Combining the held-out fluctuation with (2) and taking the finite maximum proves (1) in the statement. The zero-denominator conventions are irrelevant with probability tending to one by the cell-count bound.

Lemma~\(\mathrm{lem{:}treatment\text{-}aware\text{-}sharp\text{-}binary\text{-}elbow}\) gives uniformly along the triangular sequence
\[
 R_{d_{\mathrm e}}^A(P_{\chi,\epsilon_m})
 -R_{d_{\mathrm p}}^A(P_{\chi,\epsilon_m})
 >\frac{1389}{5120}. \tag{3}
\]
On the event that the maximum error in (1) is less than half the right-hand side of (3), the empirical minimizer is uniquely \(d_{\mathrm p}\). The probability of this event tends to one, proving (2).

On \(\{\widehat d_m^A=d_{\mathrm p}\}\), each random ratio in (3) of the statement equals its corresponding center population ratio. The treatment-aware sharp-elbow lemma identifies these with the original ratios, and its lower bound gives
\[
 \rho_{\mathrm{pred}}^A(P_{\chi,\epsilon_m})
 =\rho_{\mathrm{pred}}^{A,\mathrm{IKM}}(P_{\chi,\epsilon_m})
 >\frac{\epsilon_m^{-2}+14}{726}\longrightarrow\infty. \tag{4}
\]
Since the event on which (4) is the achieved ratio has probability tending to one, both achieved ratios diverge in probability. This proves (3) and the theorem.